\section{Technical Indicators}

\subsection{Relative Strength Index (RSI)}

The RSI indicator moves between zero and 100, plotting recent price gains versus recent
price losses. The RSI levels therefore help in estimating momentum and trend strength:
\begin{equation}
  RSI_t(n) = \frac{\sum_{k=0}^{n-1}(P_{t-i}-P_{t-i-1})\,\mathbb{1}\{P_{t-i}>P_{t-i-1}\}}{\lvert P_{t-1}-P_{t-i-1}\rvert}\times 100,
\end{equation}
where $RSI_t$ is the Relative Strength Index at time $t$, $P_t$ is the value of the index
at time $t$, and $n$ is the number of RSI periods. The RSI ranges from 0 to 100. A stock
is considered \emph{overbought} when its RSI is above 70, while it is regarded as
\emph{oversold} when the RSI is below 30.

\subsection{Moving Average Convergence Divergence (MACD)}

The MACD helps traders see the trend direction as well as the momentum of that trend. It
also provides a number of trade signals. When the MACD is above zero, the price is in an
upward phase; if the MACD is below zero, it has entered a bearish period. The indicator is
composed of two lines: the MACD line and a signal line, which moves slower. When the MACD
crosses below the signal line, it indicates that the price is falling; when the MACD line
crosses above the signal line, the price is rising. The MACD is constructed based on
moving averages: it is calculated by subtracting the longer exponential moving average
(EMA) from the shorter EMA. The EMA is defined as
\begin{equation}
  EMA_t = \frac{2}{n}\times(p_t - EMA_{t-1}) + EMA_{t-1},
\end{equation}
where $EMA_t$ is the exponential moving average at time $t$ and $n$ is the number of
periods for the EMA.

\subsection{On-Balance Volume (OBV)}

The on-balance volume indicator measures the positive and negative flow of volume over
time. It is calculated using the closing price and the trading volume of a stock.
Mathematically, the OBV of a stock at time $t$ is defined as
\begin{equation}
  OBV_t =
  \begin{cases}
    OBV_{t-1} + V_t, & C_t > C_{t-1},\\[2pt]
    OBV_{t-1} - V_t, & C_t < C_{t-1},
  \end{cases}
\end{equation}
where $C_t$ and $V_t$ are the closing price and the trading volume, respectively, at time
$t$. By definition, the OBV increases when prices rise and vice versa.

\subsection{Aroon Indicator}

The Aroon indicator measures whether the price is hitting new highs or lows over the
calculation period (typically 25). The formula for the Aroon oscillator is
\begin{align}
  Aroon\ Oscillator &= Aroon\ Up - Aroon\ Down,\\
  Aroon\ Up &= 100\times\frac{25 - \text{periods since 25-period high}}{25},\\
  Aroon\ Down &= 100\times\frac{25 - \text{periods since 25-period low}}{25}.
\end{align}
When the Aroon Up crosses above the Aroon Down, that is the first sign of a possible trend
change. If the Aroon Up hits 100 and stays relatively close to that level while the Aroon
Down stays near zero, that is positive confirmation of an uptrend.

\subsection{Stochastic Oscillator}

The stochastic oscillator measures the current price relative to the price range over a
number of periods. Plotted between zero and 100, the idea is that when the trend is up,
the price should be making new highs; in a downtrend, the price tends to make new lows.
The stochastic tracks whether this is happening:
\begin{align}
  CL_t &= P_t - \min(P_t, P_{t-1}, \dots, P_{t-s}),\\
  HL_t &= \max(P_t, P_{t-1}, \dots, P_{t-8}) - \min(P_t, P_{t-1}, \dots, P_{t-8}),\\
  RSV  &= \frac{CL_t}{HL_t}\times 100,\\
  K_t  &= \frac{2}{3}K_{t-1} + \frac{1}{3}RSV_t,\\
  D_t  &= \frac{2}{3}D_{t-1} + \frac{1}{3}K_t,
\end{align}
where $CL_t$ is measured as the lowest closing price in $N$ recent days subtracted from the
latest closing price; $HL_t$ refers to the difference between the highest and the lowest
closing price within $N$ days; $RSV_t$ is the $CL_t$ over the $HL_t$; the $K$ value is the
sum of $1/3$ of the $RSV$ value and $2/3$ of the $K$ value at lag one; and the $D$ value is
the sum of $1/3$ of the $K$ value and $2/3$ of the $D$ value at lag one. According to the
above equations, the $K$ and $D$ values range from 0 to 100. Overbought signals are emitted
as $K>80$ and $K>D$; oversold signals are emitted as $K<20$ and $K<D$.

\subsection{Comments}

\begin{figure}[H]
  \centering
  \includegraphics[width=0.78\textwidth]{z108_1.png}
  \caption{Selected technical indicators -- comparison of revenues, buy and hold vs.\ agent.}
\end{figure}

\begin{figure}[H]
  \centering
  \includegraphics[width=0.78\textwidth]{z109_1.png}
  \caption{Selected technical indicators -- actions chosen by the trained agents over time.}
\end{figure}

Comparing the graphs on the ratio of revenues between buy and hold and the trained agent,
it is obvious that the revenues are similarly distributed, so there is no advantage for
one of the strategies. However, the actions chosen over time give an insight into the fact
that the trained agents suggest clearly different decisions. It can therefore be assumed
that the trained agent applied to this share could lead to inconsistent behavior.
Comparing those results with the agent trained with all indicators from above, it is well
evident that in both described categories no better result can be achieved by only using
the 5 indicators I selected.
